%% kapa: exam=1 round=19 year=2008 period=2 form=A subject=economics subject_ko=경제원론 seq=34 number=34
%% source: Q-Net 감정평가사 기출문제 (cst00309) artlSeq=1001962 file=extracted/1008/1001962/2008년 제19회 감정평가사 1차시험 기출문제/2교시.hwp
%% conversion: kordoc 4.24.51 -> tools/gen1.py (auto); review=proofread-sonnet; flags=none
\begin{question}{34}
필립스곡선이 다음과 같고, 현재 경제는 자연실업률 수준에 있다고 한다. 예상인플레이션율($\pi^{ e }$)이 0\%인 경우와 3\%인 경우, 실업률($u$)을 자연실업률 수준보다 1\%포인트 더 낮추려면 각각 얼마만큼의 인플레이션율($\pi$)을 감수해야 하는가?
\begin{center}$\pi = 5.0 - 2u + \pi^{e}$\end{center}
\begin{choices}
  \item[] \hspace*{3.5em}0\% 인플레이션율 예상 \hspace{1.5em} 3\% 인플레이션율 예상
  \item \hspace*{6em}1.0\% \hspace*{8em} 3.0\%
  \item \hspace*{6em}2.0\% \hspace*{8em} 3.0\%
  \item \hspace*{6em}2.0\% \hspace*{8em} 5.0\%
  \item \hspace*{6em}3.0\% \hspace*{8em} 5.0\%
  \item \hspace*{6em}3.0\% \hspace*{8em} 7.0\%
\end{choices}
\end{question}
\answer{3}
